% ARCHIVED DRAFT: independently audited by root and incorporated on 2026-11-19 UTC.
% The assembled paper is authoritative; this file retains drafting provenance.
\begin{saresult}{B.077}{Proved}{Every conductor cut in a real-valued valuation domain}
Let $V$ be a characteristic-$p$ valuation domain with infinite residue
field $l$ and a nontrivial valuation
$v:K^\times\to\Gamma\subset\mathbb R$, where
$K=\operatorname{Frac}(V)$. Suppose $V$ contains $k=\mathbb F_q$.
Let $J$ be any nonzero proper ideal of $V$ and put $R=k+J$.
For every unital commutative inclusion $R\subset S$, with arbitrary
base torsion, the following are equivalent:
\[
\begin{gathered}
 \text{the polynomial co-operations exist uniquely};\\
 \theta_2\text{ is injective};\qquad (R:_S J)=R;\\
 \operatorname{Ann}_{S/R}(J)=0;\qquad \Omega(R,S)=R[X].    \tag{B77.1}
\end{gathered}
\]
All positive cases have every multivariable polynomial tensor
isomorphism. The ideal $J$ may be an open or irrational valuation cut;
it need not have a generator or a least element of its value set.
In particular this settles the previously listed Hahn example with
residue $\mathbb F_2(u)$, value group $\mathbb Q$, and $J=\{v>\sqrt2\}$.
\end{saresult}
\begin{proof}
The ring $R$ is local with maximal ideal $J$ and residue field $k$.
Let $\alpha=\inf v(J\setminus\{0\})\ge0$.
If $\alpha=0$, then $J$ is the maximal ideal of $V$. For $r\in J$ and
any $n$, choose $y\in J$ with $nv(y)<v(r)$. The quotient $r/y^n$ has
positive valuation and belongs to $J\subset R$. Thus $J$ is generated
as an $R$-ideal by its $n$-th powers for every $n$. B.069, including
its determinant calculation, gives exactly (B77.1).
We henceforth assume $\alpha>0$.

The determinant ideals are cofinal with powers of $J$.
Residue multiplicities give $D_d(R)\subset J^n$ for sufficiently
large $d$, as in B.075. Conversely choose $c\in J$ and a grid
$cv_i$ whose $v_i\in V$ have distinct residues in $l$.
Its determinant has valuation $h v(c)$, where $h=d(d+1)/2$.
Choose $n$ with $n\alpha>hv(c)+\alpha$. Every element of $J^n$ divided
by that determinant then has valuation greater than $\alpha$, hence
belongs to $J\subset R$. Thus $J^n\subset D_d(R)$.
Consequently the actual coefficient ring is
\[
                    C=\bigcup_{n\ge0}(R:_S J^n).          \tag{B77.2}
\]
The colon identity already proves the positive implication from
$(R:_S J)=R$. Assume tensor injectivity and hence epicness of $R\to C$.

Rationalization maps $C$ to an epimorphic overring $L\subset K$, with
kernel killed elementwise by powers of $J$. If $L\subset V$, the common
ideal $J$ reduces this to a field epimorphism $k\to L/J$, so $L=R$.
If $L$ contains an element of negative valuation, a power of its inverse
has valuation greater than $\alpha$ and belongs to $J$; multiplying
$V$ by that element puts it in $J$, so $V\subset L$.
The real-valued rank-one property then gives $L=K$.
If $L=R$, the ring retraction forces $C=R$.
If $L=K$, choose any nonzero $x\in J$ and apply the polynomial
idempotent construction of B.075 to a lift of $x^{-1}$.
The same nilpotent-quotient argument gives
\[
                  C\simeq K\times R/I,\qquad J^M\subset I
                    \quad\text{for some }M\ge1,           \tag{B77.3}
\]
allowing $I=R$.

Put $P(X)=X^q-X$. The monomials $X^iP(X)^j$, $0\le i<q$, $j\ge0$,
form a monic $K$-basis of $K[X]$.
Their precise $R$-valued coefficient modules are
\[
 \Omega(R,K)=\bigoplus_{0\le i<q}RX^i
 \ \oplus\!\bigoplus_{0\le i<q,\ j\ge1}
                          (J:_KJ^j)X^iP(X)^j.            \tag{B77.4}
\]
Here is a direct proof of necessity, including the valuation boundary.
Subtract the degree-less-than-$q$ interpolation polynomial of the
values on $k$; its coefficients lie in $R$, leaving all $j=0$
coefficients zero. Fix $\gamma\in v(J)$ and choose $c\in J$ of that
value. In evaluating at $b+cz$, with $b\in k$ and $z\in V$, the minimum
term valuation is
\[
                   \delta=\min_{0\le i<q,\ j\ge1}
                                 \bigl(v(a_{ij})+j\gamma\bigr).
\]
Divide by a coefficient attaining this minimum after multiplication
by $c^j$. In residue, interpolation over $b\in k$ shows that for some
$b$ the result is a nonzero polynomial in $\bar z\in l$, with no
constant term. If $\delta<0$, a value avoiding its finitely many roots
contradicts $V$-valuedness. If $\delta=0$, the unnormalized residue polynomial takes the
infinite field $l$ into $k$. Equivalently, the normalized polynomial
takes values in a finite scalar multiple of $k$. Either way a polynomial
with finite image on the infinite field $l$ is constant, again impossible. If $\delta>0$ but is outside the value set
of $J$, a value with valuation exactly $\delta$ contradicts
$R\cap\mathfrak m_V=J$. Thus $\delta$ belongs to the value set of $J$.
Every coefficient consequently satisfies $a_{ij}c^j\in J$ for every
$c\in J$. In a valuation ring, products of $j$ elements of $J$ have
valuation at least $j$ times the least valuation among those elements.
This proves $a_{ij}J^j\subset J$. Sufficiency of these conditions is
termwise evaluation, proving (B77.4).

In particular $(J:_KJ)=V$. A coefficient of negative valuation cannot
stabilize $J$, since an element with valuation sufficiently close to
$\alpha$ is sent below $\alpha$; every coefficient in $V$ does stabilize
the ideal. Therefore the $j=1,i=0$ summand of (B77.4) is $VP$.
If $I=R$ in (B77.3), choose $\beta\in V$ with residue outside $k$.
The tensor $\beta P\otimes P-P\otimes\beta P$ has zero polynomial
image and maps to a nonzero tensor in $l\otimes_k l$ under this split
coefficient projection and residue. Thus $C=K$ is excluded.

Suppose $I$ is proper. Define explicitly
\[
 H=\{F\in K[X]:F(R)\subset I\},\qquad E=R[X]+H.
\]
The latter is exactly the first projection of $\Omega(R,K\times R/I)$.
Let $H_0$ consist of the polynomials in $H$ vanishing on $k$; subtraction
of interpolation on $k$, with coefficients in $I$, maps $H$ into $H_0$.
For a nonzero $F=\sum_{0\le i<q,\ j\ge1}a_{ij}X^iP^j\in H_0$, define its real
weighted value by
\[
             w(F)=\min_{0\le i<q,\ j\ge1}\bigl(v(a_{ij})+j\alpha\bigr).
\]
In this formula $i$ ranges over $0\le i<q$; only $j$ is required to
be positive. Taking $\gamma$ down to $\alpha$ in the coefficient
criterion (B77.4) shows $w(F)\ge\alpha$.
The module $H_0$ is nonzero: if $0\ne c\in J^{M-1}$, then $cP\in H_0$.
Hence
\[
             \lambda=\inf\{w(F):0\ne F\in H_0\}
                   \quad\text{is finite and }\lambda\ge\alpha>0.
                                                               \tag{B77.5}
\]

We show $H\not\subset JE$. Choose $G\in H_0$ with
$w(G)<\lambda+\alpha/4$ and let $D=\deg G\ge1$.
For some $b_0\in k$, the ordinary polynomial $G(b_0+Y)$ has Gauss
value $w(G)$ when $Y$ is assigned weight $\alpha$.
Indeed every term of its expansion has at least that weight. Choose
a $P$-index $j$ attaining the minimum; interpolation on $k$ makes the
coefficient polynomial of that $P^j$ have minimum valuation at some
$b_0$. Its contribution to the coefficient of $Y^j$ has weight
$w(G)$, while contributions from smaller $P$-indices have strictly
larger weight because $\alpha>0$. This proves the assertion.

Choose an integer $N$ with $(N+1)\alpha>\lambda+\alpha$, and choose
$c_N\in K$ such that
\[
        0\le\varepsilon=v(c_N)+(N-1)\alpha<\frac\alpha{4D}.
                                                               \tag{B77.6}
\]
If $\alpha\in\Gamma$, choose equality zero. Otherwise $\Gamma$ is
dense in $\mathbb R$ and an element of the indicated interval exists.
For completeness, an additive subgroup of $\mathbb R$ which is not
dense has a least positive element and is cyclic; every nonzero ideal
cut in that cyclic value group is attained, so its infimum belongs
to the group. Thus these two choices cover every case.
The polynomial
\[
                       \psi_N(X)=b_0+c_NP(X)^N
\]
maps $R$ into $R$. For $P(r)\ne0$, its valuation lies in the value set
of $J$, so the nonconstant term has valuation
$\alpha+\varepsilon+N(v(P(r))-\alpha)$, in that same value set.
If the cut is open this is strictly greater than $\alpha$; if it is
closed equality is allowed.

Put $F=G\circ\psi_N$. It belongs to $H_0$, and its expansion is a
polynomial in $P^N$ with zero constant term. Its weighted value is
exactly the Gauss value of $G(b_0+Y)$ with weight
$\alpha+\varepsilon$, since distinct powers of $Y$ become distinct
powers of $P^N$. Consequently
\[
                  w(F)\le w(G)+D\varepsilon
                                  <\lambda+\alpha/2,      \tag{B77.7}
\]
and every nonzero $P$-index of $F$ is at least $N$.

If $H\subset JE$, write $F=A+\sum_i r_iF_i$, with
$A\in JR[X]$, $r_i\in J$, $F_i\in H$.
Remove the interpolation polynomials of the $F_i$ and absorb them
into $A$, so every $F_i\in H_0$.
Every coefficient of $A$ in the monic $X^iP^j$ basis belongs to $J$,
so its weight at a degree $j\ge N$ is at least
$(j+1)\alpha>\lambda+\alpha$.
At a monomial attaining $w(F)$, a summand $r_iF_i$ must therefore
have weight at most $w(F)$. This implies
\[
                   w(F_i)\le w(F)-v(r_i)
                         \le w(F)-\alpha<\lambda,
\]
contrary to (B77.5). Thus $H\not\subset JE$.

Choose $F\in H\setminus JE$. The polynomial $f=(F,0)$ survives modulo
$J\Omega$ by its first projection. The null polynomial
$g=(0,P^M\bmod I)$ survives there by its monic second projection over
$k$. Their tensor is nonzero modulo $J$ and has zero polynomial image,
excluding every proper $I$ in (B77.3).
We have proved $C=R$ under tensor injectivity. The iterated-colon
identity and B.064 now give every assertion in (B77.1).
\end{proof}
